% ============================================================
% v2 — Capitolo 2: Il progetto — la piattaforma sperimentale
% Blocco II di IV. Budget: ~7 pagine.
% Fonti verificate da codice (invariate rispetto alla v1):
%   carla_multi_agent_env.py, route_planner.py,
%   curriculum_batch_manager.py, centralized_critic.py,
%   train_mappo.yaml, curriculum_batch.yaml, levels.yaml.
% ============================================================
\chapter{Il progetto --- la piattaforma sperimentale}
\label{ch:piattaforma}

Questo capitolo descrive la piattaforma su cui poggia l'esperimento: un
ambiente di apprendimento per rinforzo multi-agente costruito su CARLA,
in cui tre veicoli e tre pedoni sono addestrati congiuntamente con l'algoritmo
MAPPO sotto due regimi alternativi --- curriculum e mixed-batch --- che
condividono ogni altro componente dello stack. L'obiettivo di design è
che il regime sia l'\emph{unica} variabile indipendente
dell'esperimento: tutto ciò che è documentato qui è comune a entrambi i
bracci, e la \cref{sec:regimi} isola il punto esatto in cui divergono.

% ------------------------------------------------------------
\section{Panoramica dell'architettura}
\label{sec:architettura}

La piattaforma consiste di quattro strati (\cref{fig:architettura}).
Il \textbf{simulatore} è CARLA 0.9.16, che gira come server autonomo
pilotato in modalità sincrona con passo di fisica fisso di \SI{0.05}{s}
(\SI{20}{Hz}): avanza solo quando l'ambiente emette un tick, il che
rende la raccolta di esperienza deterministica rispetto allo scheduling.
Il traffico scriptato non apprendente --- 15 veicoli e 30 pedoni NPC in
ogni livello --- è gestito dal Traffic Manager, anch'esso sincrono.
L'\textbf{ambiente multi-agente} (\texttt{CarlaMultiAgentEnv}) espone la
simulazione a RLlib come interfaccia multi-agente con sei agenti
apprendenti: a ogni reset spawna gli agenti, assegna a ciascuno un
percorso individuale (\cref{sec:percorsi}) e configura lo scenario
secondo il livello selezionato dal manager curriculum/batch
(\cref{sec:regimi}); un episodio dura al più 1000 passi di controllo,
cioè \SI{50}{s} di tempo simulato.

Lo \textbf{stack di apprendimento} è MAPPO su Ray/RLlib 2.10.0 con
PyTorch 2.7. Vengono apprese due policy, una condivisa dai tre veicoli e
una dai tre pedoni --- condivisione di parametri entro ciascuna classe,
nessuna condivisione tra classi, poiché le due hanno spazi di
osservazione e azione diversi. Ogni actor è un percettrone multistrato
con due strati nascosti da 256 unità; un critic centralizzato riceve
un'osservazione globale a 225 dimensioni che concatena informazione per
agente (\cref{sec:osservazioni}), con encoder MLP di default e una
variante a grafo come asse sperimentale secondario (\cref{sec:critic}).
Infine il \textbf{logging}: alla fine di ogni episodio il processo di
addestramento accoda un record JSON \emph{per agente} --- sei per
episodio --- a un file \texttt{episodes.jsonl} contenente ragione di
terminazione, completamento ed efficienza del percorso, velocità e
diagnostica di generazione del percorso. Tutti i risultati del
\cref{ch:risultati} sono calcolati da questi record.

\begin{figure}[t]
  \centering
  \begin{tikzpicture}[
      node distance=5mm and 10mm,
      box/.style={draw, rounded corners=2pt, align=center,
                  minimum width=48mm, minimum height=9mm, font=\small},
      lbl/.style={font=\scriptsize, midway},
      >={Stealth[length=2.5mm]}]
    \node[box] (carla) {Server CARLA 0.9.16\\ \scriptsize Town03 (train) / Town05 (test), TM};
    \node[box, below=of carla] (env) {\texttt{CarlaMultiAgentEnv}\\ \scriptsize 3 veicoli + 3 walker, route planner};
    \node[box, right=of env] (mgr) {Manager curriculum\\ / batch \scriptsize (livello per episodio)};
    \node[box, below=of env] (rllib) {Ray/RLlib 2.10 --- MAPPO (CTDE)\\ \scriptsize policy veicolo, policy pedone, critic centralizzato};
    \node[box, below=of rllib] (log) {\texttt{episodes.jsonl}\\ \scriptsize 6 record/episodio};
    \node[box, right=of log] (eval) {Pipeline di valutazione\\ \scriptsize $4$ scenari $\times$ 100 ep};
    \draw[->] (env) -- node[lbl, right] {tick / controllo} (carla);
    \draw[<-] ([xshift=-8mm]env.north) -- node[lbl, left] {sensori, stato} ([xshift=-8mm]carla.south);
    \draw[->] (mgr) -- node[lbl, above] {livello} (env);
    \draw[->] (env) -- node[lbl, right] {oss., reward} (rllib);
    \draw[<-] ([xshift=-8mm]env.south) -- node[lbl, left] {azioni} ([xshift=-8mm]rllib.north);
    \draw[->] (rllib) -- (log);
    \draw[->] (rllib) -| node[lbl, above, pos=0.25] {checkpoint} (eval);
  \end{tikzpicture}
  \caption{Architettura a componenti. Il manager curriculum/batch è
  l'unico componente che differisce tra i due bracci sperimentali, e
  solo nella regola di selezione del livello (\cref{sec:regimi}).}
  \label{fig:architettura}
\end{figure}

% ------------------------------------------------------------
\section{Un mondo condiviso: la copertura parallela dei task}
\label{sec:mondo}

I sei agenti apprendenti non cooperano verso un obiettivo congiunto: a
ciascuno è assegnato il \emph{proprio} percorso ed è valutato
individualmente sul suo completamento. Poiché tutti e sei i percorsi
sono collocati nello stesso mondo vivo, un singolo episodio esercita
simultaneamente sei micro-task diversi --- geometrie stradali diverse,
interazioni diverse con il traffico scriptato e, soprattutto,
interazioni reciproche tra agenti apprendenti, per esempio un veicolo
che dà la precedenza a un pedone anch'esso in addestramento. Chiamiamo
questo design \emph{copertura parallela dei task}: per episodio di
orologio il sistema raccoglie esperienza su sei istanze di task anziché
una. Rispetto ai benchmark cooperativi che motivano MAPPO, il setting è
meglio descritto come \emph{coesistenza} di task indipendenti in un
mondo condiviso --- nessuna reward di squadra, nessuna competizione
esplicita, ma un accoppiamento reale attraverso lo stato del mondo.

Ne seguono due avvertenze. Primo, i sei campioni per episodio sono
\emph{correlati}, perché gli agenti condividono lo stato del mondo: per
questo il \cref{ch:metodo} non usa mai i conteggi di record come se
fossero osservazioni indipendenti. Secondo, dal punto di vista di ogni
singolo agente l'ambiente è \emph{non-stazionario} durante
l'addestramento, perché le policy degli altri cinque cambiano nel tempo
--- la motivazione standard del critic CTDE. La copertura parallela
resta comunque una proprietà della piattaforma, identica in entrambi i
regimi: parte della cornice sperimentale, non una variabile.

\paragraph{Percorsi per agente e seeding.}
A ogni reset ciascun agente riceve un percorso indipendente la cui
lunghezza target è fissata dal livello attivo. La generazione ha un seed
deterministico per agente con una sequenza costruita da:

$\langle\text{seed traffico},\ \text{contatore di reset},\
\mathrm{crc32}(\text{id agente})\rangle$

cosicché due run lanciate con lo stesso seed sullo stesso codice vedano 
la stessa sequenza di estrazioni e le estrazioni siano disaccoppiate tra 
agenti. È questa proprietà a rendere i confronti A/B appaiati \emph{per costruzione}, 
e su di essa il protocollo del \cref{ch:metodo} fa affidamento.

\paragraph{Tassonomia delle terminazioni.}
Ogni agente termina il proprio episodio in modo indipendente con
esattamente una di sei ragioni: \texttt{route\_complete} (l'\emph{unico}
esito contato come successo); \texttt{route\_short}, cioè la fine di un
percorso risultato più corto del target oltre una tolleranza --- un
fallback degenere che \emph{non} conta come successo, così da prevenire
l'inflazione della metrica tramite percorsi banalmente brevi;
\texttt{collision}; \texttt{offroad} (il veicolo si è spostato più di
\SI{5}{m} dal centro della corsia più vicina); \texttt{stuck} (mancanza
prolungata di progresso); e \texttt{timeout}. Ogni episodio produce
quindi esattamente sei record a livello di agente.

% ------------------------------------------------------------
\section{Spazi di osservazione e azione}
\label{sec:osservazioni}

Ogni agente riceve un vettore di osservazione compatto e specifico per
classe, con tutte le componenti normalizzate in $[-1,1]$ e sanificate
contro valori non finiti: 47 dimensioni per i veicoli, 26 per i pedoni
(\cref{tab:obs}). Il design architetturale è strutturato come segue:
cinematica dell'ego e dei vicini più prossimi, in cui entrambe le
classi osservano una breve \emph{anteprima} del proprio percorso ---
offset relativi dei waypoint successivi nel riferimento dell'ego ---
insieme all'errore di heading e, per i veicoli, all'angolo della curva
imminente: la geometria locale del task assegnato, senza nulla degli
obiettivi altrui. In secondo luogo, una routine condivisa stima lungo il
corridoio frontale un rischio di time-to-collision e una frazione di
occupazione, calcolati separatamente contro i veicoli (corridoio
\SI{3.5}{m}, \SI{40}{m} avanti, orizzonte \SI{5}{s}) e contro i pedoni
(\SI{3.0}{m}, \SI{30}{m}, \SI{4}{s}); gli stessi quattro valori lato
veicolo sono riusati dalla reward come gate di sicurezza
(\cref{sec:reward}), così che la quantità che modula la reward sia
\emph{sempre} visibile alla policy. Infine, le ultime tre dimensioni del
veicolo espongono stato interno dell'ambiente da cui la reward dipende
--- passi senza progresso, penalità anti-loop e tempo rimanente ---
senza le quali le penalità agirebbero su informazione che l'agente non
può osservare, un difetto di aliasing di stato identificato e corretto
durante lo sviluppo (\cref{sec:gate}).

\paragraph{Spazi di azione.}
Entrambe le classi usano azioni continue bidimensionali. L'azione del
veicolo è $[a_{\mathrm{tb}}, a_{\mathrm{steer}}] \in [-1,1]^2$, dove
$a_{\mathrm{tb}}$ positivo mappa sull'acceleratore e negativo sul freno;
quella del pedone è $[a_{\mathrm{speed}}, a_{\Delta\mathrm{h}}]$, con
$a_{\mathrm{speed}} \in [0,1]$ che scala una velocità massima di
camminata di \SI{5.0}{m/s} e $a_{\Delta\mathrm{h}} \in [-1,1]$ che
comanda un cambio di heading fino a $\pm 45^{\circ}$ per passo. Per il
controllo continuo la policy è una gaussiana diagonale: il campionamento
non è un dettaglio implementativo ma parte dell'oggetto appreso, un
punto che diventerà concreto nella \cref{sec:validita}.

\paragraph{Osservazione globale per il critic.}
Durante il solo addestramento il critic riceve una concatenazione a slot
fissi delle osservazioni di tutti gli agenti seguita da una maschera di
vitalità:
$[\,v_0\,|\,v_1\,|\,v_2\,|\,p_0\,|\,p_1\,|\,p_2\,|\,\text{alive}_{6}\,]$,
cioè $3 \times 47 + 3 \times 26 + 6 = 225$ dimensioni. La maschera è
necessaria perché gli agenti terminano in modo indipendente: gli slot di
quelli terminati vengono mascherati anziché riusati silenziosamente. Gli
actor non vedono mai questo vettore, preservando l'esecuzione
decentralizzata.

% ------------------------------------------------------------
\section{Design della reward}
\label{sec:reward}

Le reward sono per agente e identiche nei due regimi; la
\cref{tab:reward} elenca ogni componente con trigger e ampiezza, e la
forma finale è l'esito delle iterazioni guidate da gate documentate
nella \cref{sec:gate} --- nessuna pretesa di ottimalità è attribuita ai
singoli coefficienti. Per entrambe le classi il segnale dominante è un
bonus sparso per waypoint raggiunto, completato da uno shaping denso
proporzionale alla riduzione per passo della distanza dal waypoint
successivo, con i waypoint saltati penalizzati così che il progresso non
sia manipolabile con scorciatoie.

Il blocco empiricamente più delicato è lo \textbf{shaping anti-stallo
condizionato all'hazard}, che contrasta l'\emph{equilibrio difensivo} in
cui un veicolo si ferma indefinitamente per evitare qualunque rischio.
Tutti i suoi termini sono sospesi mentre il veicolo attende
legittimamente a un semaforo rosso o giallo. Lo scalare
$\mathit{hazard} = \max(\text{TTC}_{\mathrm{veh}}, \text{occ}_{\mathrm{veh}},
\text{TTC}_{\mathrm{ped}}, \text{occ}_{\mathrm{ped}})$ riassume il canale
di anteprima della \cref{sec:osservazioni}, e un booleano
$\mathit{safe\_to\_push} = (\mathit{hazard} < 0.85)$ fa da gate ai
termini di spinta: una penalità di inattività sotto 2.5~km/h; un bonus
di partenza proporzionale ad acceleratore, allineamento e a un fattore
di urgenza crescente con il tempo trascorso senza progresso; e uno
shaping di velocità minima attorno a un target di 8~km/h. Sopra la
soglia di hazard ogni incentivo di spinta scompare, quindi la policy non
è mai premiata per accelerare in un corridoio rischioso;
indipendentemente dal gate, una rampa di non-progresso e una penalità
anti-loop rendono lo stallo prolungato e il girare in tondo strettamente
non redditizi.

I pedoni, infine, guadagnano un bonus per passo mentre camminano entro
una \textbf{banda di comfort} di $[1.2, 2.6]$~m/s, intenzionalmente
ampia: a 2.0~m/s un pedone copre \SI{100}{m} --- il target del livello
hard --- entro l'orizzonte di \SI{50}{s}. Il suo design asimmetrico
(solo bonus) e la sua interazione con il gate di hazard dei veicoli sono
un caso concreto di accoppiamento tra classi, analizzato empiricamente
nel \cref{ch:risultati}.

% ------------------------------------------------------------
\section{Generazione dei percorsi e livelli di difficoltà}
\label{sec:percorsi}

La difficoltà è definita dalla \emph{lunghezza target} dei percorsi per
agente, con tutti gli altri parametri di scenario tenuti costanti: i tre
livelli di addestramento collocano i target di veicoli e pedoni a 30
(easy), 60 (medium) e \SI{100}{m} (hard), sempre su Town03 e sempre con
lo stesso traffico scriptato; un quarto livello di sola valutazione
(\emph{test}) usa Town05 con target di \SI{80}{m}.

I percorsi dei veicoli sono pianificati sulla topologia stradale con il
global route planner di CARLA: le destinazioni candidate sono campionate
attorno alla distanza target e accettate solo se la lunghezza del
cammino ottimo ricade in $[0.5, 2.0] \times$ il target, con fino a 32
tentativi per reset; se nessuna sopravvive, l'ambiente ripiega su un
generatore legacy a catena di waypoint. I percorsi dei pedoni sono
invece catene di segmenti di marciapiede connesse attraverso
attraversamenti pedonali, plasmate da due vincoli: le catene sotto
$0.5 \times$ il target vengono rigettate --- la stessa guardia che
sostiene la retrocessione \texttt{route\_short} --- e gli
attraversamenti per percorso sono limitati a tre, un tetto
all'esposizione sulla carreggiata fissato empiricamente
(\cref{sec:gate}). La fattibilità dei percorsi pedonali su Town03 è una
\emph{limitazione strutturale} della piattaforma --- la topologia dei
marciapiedi spesso non può sostenere i target nominali medium e hard ---
ed è analizzata nella \cref{sec:pedoni}.

% ------------------------------------------------------------
\section{I due regimi di addestramento}
\label{sec:regimi}

Il regime è la variabile indipendente dell'esperimento principale.
Entrambi sono implementati dallo stesso componente manager, consultato
una volta per reset per selezionare il livello; condividono le
definizioni di livello della \cref{sec:percorsi}, il budget totale di
$3 \times 10^6$ passi d'ambiente e ogni altro componente descritto in
questo capitolo. Differiscono \emph{solo} nella regola di selezione.

\paragraph{Braccio curriculum: progressione condizionata alla competenza.}
L'addestramento parte con il solo livello easy sbloccato. I livelli
superiori si sbloccano attraverso un gate di competenza valutato su una
finestra scorrevole degli ultimi 50 episodi: il tasso di successo a
finestra è calcolato separatamente per la policy veicolo e per quella
pedone, e il gate vincola sul \emph{minimo} dei due --- la progressione
richiede che la classe più debole sia competente, non la media. Medium
si sblocca a un tasso di successo a finestra di almeno il 70\%, hard al
60\%; la finestra deve essere piena, e una quota minima del budget deve
già essere stata spesa ai livelli precedenti (20\% su easy prima di
medium, 28\% su medium prima di hard). Due decisioni di design meritano
giustificazione. La prima è il gating \emph{a finestra} e non
cumulativo: il tasso cumulativo è un integratore in ritardo, trascinato
verso il basso dai fallimenti di avvio a freddo, e una policy può essere
competente \emph{adesso} mentre la sua media dalla nascita è ancora ben
sotto soglia. La seconda sono i \emph{cap di pressione}: se una soglia
non viene mai raggiunta, il livello si sblocca comunque una volta
consumata una quota fissa del budget (30\% per medium, 70\% per hard),
così che il braccio curriculum non possa fermarsi e terminalmente su easy.
Una volta sbloccati, i livelli sono campionati con pesi che normalizzano
a $0.30/0.35/0.35$, e quelli appena sbloccati passano per una breve fase
di \emph{probation} a pressione ridotta (\cref{fig:fsm}).

\paragraph{Braccio batch: alternanza stratificata.}
Il braccio mixed-batch è il \emph{controllo senza ordinamento}: tutti e
tre i livelli sono disponibili dal primo episodio, campionati per
mescolamento stratificato senza reinserimento con finestra --- ogni
finestra consecutiva di tre episodi è una permutazione uniformemente
casuale di \{easy, medium, hard\}. Ne risulta esposizione uniforme
$1/3$ con sbilanciamento massimo di un episodio per campionatore,
evitando sia la deriva del campionamento i.i.d.\ sia qualunque
ordinamento deterministico fisso.


\begin{figure}[t]
  \centering
  \begin{tikzpicture}[
      state/.style={draw, rounded corners=2pt, align=center,
                    minimum width=31mm, minimum height=11mm, font=\small},
      gate/.style={font=\scriptsize, align=center},
      >={Stealth[length=2.5mm]}, node distance=15mm]
    \node[state] (s1) {solo easy\\ \scriptsize campionamento: easy 1.0};
    \node[state, right=of s1] (s2) {easy $+$ medium\\ \scriptsize probation, poi\\ \scriptsize $\propto 0.30/0.35$};
    \node[state, right=of s2] (s3) {easy $+$ medium $+$ hard\\ \scriptsize probation (hard $\times$0.85), poi\\ \scriptsize $\propto 0.30/0.35/0.35$};
    \draw[->] (s1.north east) to[bend left=26]
      node[gate, above] {SR fin.\ $\geq 0.70$\\ $\wedge$ quota easy $\geq 0.20$} (s2.north west);
    \draw[->, dashed] (s1.south east) to[bend right=26]
      node[gate, below] {cap: 30\% budget} (s2.south west);
    \draw[->] (s2.north east) to[bend left=26]
      node[gate, above] {SR fin.\ $\geq 0.60$\\ $\wedge$ quota med.\ $\geq 0.28$} (s3.north west);
    \draw[->, dashed] (s2.south east) to[bend right=26]
      node[gate, below] {cap: 70\% budget} (s3.south west);
  \end{tikzpicture}
  \caption{Progressione del curriculum. Archi pieni: gate di competenza
  sul tasso di successo a finestra (minimo sulle due policy, finestra di
  50 episodi piena) con soglie minime di quota di budget. Archi
  tratteggiati: cap di pressione che garantiscono la progressione entro
  il budget.}
  \label{fig:fsm}
\end{figure}


Sotto questa costruzione i due bracci consumano lo stesso budget sulla
stessa distribuzione di task in attesa; ciò che differisce è
\emph{quando} ogni parte di quella distribuzione viene presentata, e se
la presentazione è condizionata alla competenza misurata. Qualunque
differenza comportamentale è quindi attribuibile alla politica di
ordinamento, a meno della varianza da run a run --- quantificata nella
\cref{sec:gate} e discussa nella \cref{sec:minacce}.

% ------------------------------------------------------------
\section{Encoder del critic, ottimizzazione e valutazione}
\label{sec:critic}

Nella configurazione di default il critic centralizzato è un MLP che
mappa l'osservazione globale a 225 dimensioni attraverso due strati
nascosti da 256 unità verso un valore scalare. L'asse sperimentale
secondario sostituisce questo encoder con uno su grafo, mentre
\emph{entrambi gli actor restano gli identici MLP} della
\cref{sec:architettura}: qualunque differenza comportamentale tra le due
varianti è quindi attribuibile alla stima del valore durante
l'addestramento e non alla capacità della policy, e a tempo di
esecuzione le due varianti sono indistinguibili. L'encoder GNN riseziona
l'osservazione globale nei sei slot per agente, proietta ciascuno
attraverso uno strato lineare per classe in un embedding di nodo a 64
dimensioni e applica due round di message passing con aggregazione a
media in stile GraphSAGE~\cite{hamilton2017graphsage}: ogni nodo combina
il proprio embedding con la media degli embedding degli altri agenti
\emph{vivi}, e un mean-pool mascherato produce l'input della testa di
valore. La maschera di vitalità guida aggregazione e pooling, così che
gli agenti terminati escano dal riassunto del critic invece di
contribuire slot stantii. La codebase fornisce anche un encoder ad
attenzione e la normalizzazione PopArt dietro flag (entrambi disabilitati 
e lasciati come finetuning futuro)

\paragraph{Ottimizzazione.}
Tutte le run della campagna condividono la stessa configurazione,
byte-identica tra i bracci salvo i flag di regime ed encoder: learning
rate $5\times10^{-4}$; $\gamma = 0.997$; GAE $\lambda = 0.95$; clip PPO
$\varepsilon = 0.2$ con penalità KL aggiuntiva; batch di addestramento
8{.}000 passi, minibatch 256, 15 epoche SGD per batch. Il coefficiente
di entropia segue uno schedule da 0.03 a 0.005 ancorato all'83\% del
budget, espresso come frazione così che run diagnostiche brevi e run
complete attraversino lo stesso profilo di esplorazione. La
configurazione completa è nell'\cref{app:config}.

\paragraph{Pipeline di valutazione.}
I risultati sono misurati a tre livelli, tutti calcolati dai record a
livello di agente. Il \emph{flusso di addestramento} fornisce le
metriche cumulative su tutti gli episodi: sono on-policy e, nel braccio
curriculum, calcolate su un mix di livelli che cambia nel tempo, quindi
caratterizzano il processo di addestramento e non l'abilità finale. Il
\emph{ricalcolo statico} è un passaggio offline che ricalcola gli
aggregati da disco, deduplica i record e verifica l'integrità. La
\emph{valutazione finale del checkpoint} ricarica il modello addestrato
e lo esegue su quattro scenari fissi da 100 episodi ciascuno --- easy,
medium e hard su Town03, più il livello test su Town05 --- ed è l'unico
confronto a distribuzione appaiata tra i bracci; gli episodi sono
valutati con la policy \emph{stocastica}, con re-seeding per episodio
dello scenario e del campionatore di azioni, perché valutare la media
deterministica collasserebbe la policy su una singola traiettoria per
scenario (\cref{sec:validita}).
